\documentclass[a4paper,10pt]{article}
\usepackage{jheppub} % for details on the use of the package, please see the JINST-author-manual
\usepackage{lineno}
\usepackage{amsmath,amsthm,amsfonts,amssymb,amscd,physics,cancel,mathtools}
\usepackage{tcolorbox}
\usepackage{marginnote,tensor}
\usepackage{tcolorbox}
%~~~~~~~~~ Document setup
% \usepackage[spanish]{babel} % English formatting
\usepackage[utf8]{inputenc} % Standard encoding
% \usepackage[a4paper,left=3cm,bottom=3cm]{geometry} % Page formatting
\usepackage{indentfirst} % Indents the first paragraph
\usepackage{amsmath} % Maths type package
\usepackage{bm} % Bold font maths
\usepackage{graphicx} % Advanced graphics package
\usepackage[export]{adjustbox} 
\usepackage{pdflscape} % Make pages landscape
\usepackage{fancyhdr} % Fancy headers
% \usepackage[colorlinks=true,citecolor=blue,urlcolor=blue,linkcolor=black]{hyperref} % Link colours
%\usepackage{natbib} % Bibliography
% \usepackage{flafter} % Reference any 'float'
% \usepackage[framemethod=tikz]{mdframed} % Box off stuff
\usepackage{color} % Colour support
\usepackage{wrapfig} % Text flowing around figures
\usepackage{lipsum} % Generates meaningless text
\usepackage{xcolor}
%\usepackage{biblatex}
%\usepackage[backend=bibtex]{biblatex}
%\addbibresource{bibliography.bib}
%\hypersetup{colorlinks=true, linkcolor=blue}

\newtheorem{ej}{Example}[section]
\newtheorem{sol}{Solution}[section]
\newtheorem{dem}{Proof}[section]
\newtheorem{prop}{Propiedad}[section]

\def\a{\alpha}
\def\b{\beta}
\def\g{\gamma}
\def\G{\Gamma}
\def\d{\delta}
%\def\D{\Delta}
%\def\e{\eta}
\def\la{\lambda}
\def\La{\Lambda}
\def\k{\kappa}
\def\m{\mu}
\def\n{\nu}
\def\r{\rho}
\def\p{\rho}
\def\o{\omega}
\def\s{\sigma}
\def\S{\Sigma}
\def\t{\tau}
\def\p{\pi}
\def\f{\phi}
\def\vf{\varphi}
\def\ep{\epsilon}
\def\th{\theta}
\def\Th{\Theta}
\def\z{\zeta}
\def\id{\mathrm{I}}
\def\M{\mathcal{M}}

\newcommand{\e}{\mathrm{e}}
\newcommand{\I}{\mathrm{I}}
\newcommand{\C}{\mathcal{C}}
\newcommand{\D}{\mathrm{D}}
\newcommand{\oo}{\mathring}
\newcommand{\oomega}{\mathring{\omega}}
\newcommand{\eabc}{\epsilon_{abc}}

%-----COLORS LIST ------
\definecolor{azure(colorwheel)}{rgb}{0.0, 0.5, 1.0}
\definecolor{DarkViolet}{RGB}{148,0,211}
\definecolor{myDarkBlue}{rgb}{0,0.1,0.7}
\definecolor{DarkBlue}{RGB}{0,0,153}
\definecolor{amber}{rgb}{1.0, 0.49, 0.0}
\definecolor{amaranth}{rgb}{0.9, 0.17, 0.31}
\definecolor{nicered}{rgb}{0.7,0.1,0.1}
\definecolor{brown}{rgb}{0.5,0.1,0.1}
\definecolor{nicegreen}{rgb}{0.0,0.3,0.0}
\definecolor{tealgreen}{rgb}{0.0, 0.51, 0.5}
\def\red#1{{\color{red} #1}}
\def\green#1{{\color{green} #1}}
\def\blue#1{{\color{blue} #1}}
\def\orange#1{{\color{orange} #1}}
%----------------------
\newcommand{\mycolor}{DarkViolet}
\def\myColor#1{{\color{\mycolor} #1}}
\definecolor{tclr}{RGB}{148,0,211}
%----------------------
\newcommand{\corr}[1]{\textcolor{nicered}{#1}}
\newcommand{\nick}[1]{\textcolor{olive}{#1}}
\newcommand{\teo}[1]{\textcolor{azure(colorwheel)}{#1}}
\newcommand{\chteo}[2]{\corr{\st{#1}} \teo{(#2)}}
\newcommand{\bako}[1]{\textcolor{DarkViolet}{#1}}
\newcommand{\than}[1]{\textcolor{magenta}{#1}}

%----------------------
\usepackage{hyperref}
\hypersetup{colorlinks,bookmarksopen,
	bookmarksnumbered,
	citecolor={nicered},
	linkcolor={myDarkBlue},
	urlcolor={blue},
	pdfstartview=FitH}


% \arxivnumber{1234.56789} % if you have one

\title{\boldmath Gravitational Instantons}

% Collaborations

%% [A] If main author
%% \collaboration{\includegraphics[height=17mm]{collabroation-logo}\\[6pt]
%%  XXX collaboration}

%% or
%% [B] If "on behalf of"
%% \collaboration[c]{on behalf of XXX collaboration}


% Authors
% The "\note" macro will give a warning: "Ignoring empty anchor...", you can safely ignore it.

%% [A] simple case: 2 authors, same institution
%% \author[1]{A. Uthor\note{Corresponding author.}}
%% \author{and A. Nother Author}
%% \affiliation{Institution,\\Address, Country}

%% or, e.g.
%% [B] more complex case: 4 authors, 3 institutions, 2 footnotes
%% \author[a,b]{F. Irst,\note{Now at another university}}
%% \author[c]{S. Econd,}
%% \author[a,2]{T. Hird\note{Also at Some University.}}
%% \author[c,2]{and Fourth}
%% \affiliation[a]{Institution_1,\\Address, Country}
%% \affiliation[b]{Institution_2,\\Address, Country}
%% \affiliation[c]{Institution_3,\\Address, Country}

\author{Borja Diez}
\affiliation{Universidad Arturo Prat}
% \affiliation{Another University,\\
% different-address, Country}

% E-mail addresses: only for the corresponding author
\emailAdd{borjadiez1014@gmail.com}

\abstract{Abstract...}




\begin{document}
\maketitle
%\tableofcontents
%\flushbottom

%\section{Differentiable manifolds}
\subsection{From Topological Spaces to Differentable Manifolds}
%TODO agregar imagen
It is assumed that the reader is acquainted with the notion of a topological space as a structure on which one can define a neighborhood and continuous functions. A \textbf{homeomorphism} between two topological spaces is a 1-1 map $\varphi: X\to Y$ for which both $\varphi$ and its inverse $\varphi^{-1}$ are continuous. If $\varphi$ and $\varphi^{-1}$ are continuously differentiable then $\varphi$ is called a \textbf{diffeomorphism}.

A $D$-dimensional manifold $M^D$ is a topological space that locally has the properties of a $D$-dimensional Euclidean space $\mathbb{R}^D$ : A neighborhood of a point in $M^D$
can continuously be mapped in a one-to-one way to the neighborhood of a point in
$\mathbb{R}^D$. To be more precise, introduce a \textbf{chart} $(U_\alpha,\varphi_\alpha)$ as a homeomorphism $\varphi_\alpha$ from
an open set $U_\alpha\subset M^D$ into an open set $R_\alpha\subset \mathbb{R}^D$. Two charts are compatible if the overlap maps are diffeomorphims $(\varphi_1\cdot \varphi_2\in C^\infty,\varphi_2\cdot \varphi_1^{-1}\in C^\infty)$ unless $U_1\cap U_2=\emptyset$. A set of compatible charts covering $M^D$ is called an atlas. In every chart the manifold can be equipped with a coordinate system: for $x\in M^D$ the coordinates are $x^\mu=\varphi(x)\in \mathbb{R}^D$. The naming makes it clear what one is aiming at. For instance the surface of a sphere, although not being homeomorphic to a plane, locally has enough smoothness to be mapped into an atlas. One chart is not sufficient since there will always be a point on the sphere that cannot be projected to the plane.

In manuscript will only treat finite-dimensional manifolds. One possibility of extending the notion of manifolds to infinite dimensions is to consider Banach manifolds modeled on Banach spaces. It is also assumed that we are dealing with $C^\infty$ manifolds. In certain contexts it might suffice that the charts are $C^k$-related. Also complex manifolds are investigated in mathematics and applied to modern theoretical physics (catchword: Kähler manifolds). In these the transition functions are required to be analytic.

\subsection{Tensor Bundles}
On a manifold one can erect tensor bundles as “superstructures” by starting with defining the tangent and cotangent spaces of a manifold.

\subsubsection{Tangent Bundle and Vector Fields}
We are interested in the notion of vectors on a manifold $M$ (henceforth I will mostly drop the index for the dimension of the manifold and for the Euclidean space). The idea is to introduce these as tangent vectors of curves 'through' $x\in M$: A curve through a point $x$ is a smooth mapping of an interval $I=[0,1]\subset\mathbb{R}$ to the manifold:
\begin{equation}
  C=\mathbb{I}\to M\qquad t\mapsto C(t)\qquad \mbox{with}\qquad C(0)=x
\end{equation}
The coordinates of this curve are $x^\mu(C(t))$, and the tangent vector to this curve is
\begin{equation}
  \dv{t}x^\mu(C(t))
\end{equation}
Since one can have more then one curve with $C(0)=x$, the proper definition is: A \textit{tangent vector} $x\in M$ is an equivalence class of curves in $M$, where the equivalence relation between two curves is that they are tangent at the point $x$. Another-equivalent- definition is to understand a tangent vector as a directional derivative: Consider functions $f\in \mathcal{F}M$, that is $f:M\to \mathbb{R}$. The change of $f$ along a curve is given by
\begin{equation}
  \dv{t}f(C(t)),\qquad \mbox{locally}\qquad \pdv{x^\mu}f\dv{x^\mu(C(t))}{t}
\end{equation}

In defining
\begin{equation}
  X=(X^\mu\partial_\mu)\qquad \mbox{with}\qquad X^\mu =\dv{x^\mu(C(t))}{t}
\end{equation}
we can write $\dv{t}f(C(t))=Xf$. For every point along the curve we take this expression to define the differential operator $X_x$ as the tangent vector to the manifold in $x\in M$. All tangent vectors at a point in the manifold can be shown to build a vector space $\mathfrak{X}_xM$ isomorphic to $\mathbb{R}^D$. The natural basis in $\mathfrak{X}_xM$ is the coordinate or \textbf{holonomic} basis $\{\partial_\mu\}$. But of course any other (\textbf{anholonomic}) basis $\{e_I}
\}$





























\section{Yang-Mills Instantons}
\subsection{Eguchi-Gilkey-Hanson}
From a geometric point of view, the path integral has the advantage of being able to take the global topology of the gauge potentials into account, while the canonical perturbation theory approach to quantization is sensitive only to the local topology. At present, a mathematically precise theory of path integration can be formulated only for spacetimes with positive signatures. We refere to those spacetimes as \textit{Euclidean} manifolds. Physically meaningful answers are obtainable by continuing the results of the Euclidean path integration back to the Minkowski regim.

In the Euclidean path-integral formulation approach to quantization, each field configuration $\f(x)$ is weighted by the \textit{Boltzmann factor} $e^{-S[\f ]}$.

For Yang-Mills theories, the Euclidean action is
\begin{equation}
  S[A_\m ]=\frac{1}{4}\int_{\mathcal{M}}\dd^4x\sqrt{g}\Tr(F_{\m\n }F^{\m\n })
\end{equation}
which is positive definite. The contribution of each gauge potential or connection $A_\m (x)$ to the path integral is therefore bounded and well-behaved.

The complete generating functional for the transition amplitudes of a theory is obtained by symming over all inequivalent configurations. Since the first-order functional variation of the action vanishes for solutions of the equations of motion, these configurations correspond to stationary points in the functional space. Therefore, in the path-integral approach, we first seek solutions to the Euclidean field equations with minimum action and the compute quantum-mechanical fluctuations around them.




In order to find the minimum action configurations of the Yang-Mills theory, let us consider the inequality
\begin{equation}
  \int_{\mathcal{M}}\dd^4x\sqrt{g}\Tr(F_{\m\n }-\tilde{F}_{\m\n })^2\geq 0
\end{equation}
This bound is satured by the self-dual field configurations
\begin{equation}
  F_{\m\n }=\pm \tilde{F}_{\m\n }
\end{equation}
In fact, these field configurations solve the Yang-Mills field equations since the Bianchi identities imply the field equations. The action now becomes
\begin{equation}
  S=-\frac{1}{2}\int \Tr F\wedge *F=\mp \int\Tr F\wedge F=4\p |k|
\end{equation}
where
\begin{equation}
  -C_2=k=-\frac{1}{8\p }\Tr F\wedge F
\end{equation}
is the integral of the second Chern class. Hooft [1976a] called such special field configurations \textit{instantons} since in the case $|k|=1$ their field strenght is centered around some point in space-time and thus attains its maximum value at some \textit{instant of time}.

\subsubsection{Physical interpretation of instantons}
The instanton can be interpreted as a quantum-mechanical tunnelin phenomenon in Yang-Mills gauge theories. It induces a transition between homotopically inequivalent vacua. The true ground state of Yang-Mills theory then becomes a coherent mixture of all these vacuum states.


\subsection{Instantons in Gauge Theories}
Correlators in Quantum Field Theories are described by path integrals over all possible field configurations
\begin{equation}
  \ev{\prod_i\mathcal{O}(x_i,t_i)}=\int\mathcal{D}\f \prod_i\mathcal{O}(x_i,t_i)e^{\frac{i}{\hbar}S(\f )}
\end{equation}
For a gauge theory
\begin{equation}
  S=-\frac{1}{4g^2}\int\dd^4xF^{a}_{\m\n }F^{a\m\n }+\cdots
\end{equation}
In the classical limit $\hbar\to 0$, the integral is dominated by the saddle point of the action $\d S=0$. To compute the contribution of the saddle point to the integral is convenient to perform the analytic continuation
\begin{equation}
  t_E=it
\end{equation}
and write
\begin{equation}
  iS=-\frac{1}{4g^2}\int\dd t_E\dd^3xF^{a}_{mn}F^{a}_{mn}+\cdots=-S_E
\end{equation}
With respect to $S(\f )$, $S_E(\f )$ has the advantage of being positive defined. The classical limit of the path integral
\begin{equation}
  \ev{\prod_i\mathcal{O}(x_i,-it_{Ei})}=\int\mathcal{D}\f \prod_i\mathcal{O}(x_i,-it_{Ri})e^{-\frac{1}{\hbar}S(\f )}
\end{equation}
is then dominated by the minima of $S_E(\f )$. A solution of the Euclidean equations of motion is called an \textit{instanton}.


Besides its applications in the computation of path integrals , instantons can be
used also to compute tunnelling effects between different vacua of a quantum field
theory. At low energies, the energy levels of a particle moving on a potential $V(x)$
can be approximated by those of the harmonic oscillator for a particle moving on
a quadratic potential $V(x)\approx \frac{\omega_0}{2}(x-x_0)^2$ around a minima at $x=x_0$. In presence of a tunnelling between two vacua, each harmonic oscillator energy levels split into two with energies
\begin{equation}
  E_n=\left(n+\frac{1}{2}\pm\frac{1}{2}\Delta_n\right)\omega_0,\qquad \Delta_n\sim e^{-S_E}
\end{equation}
with $S_E$ the Euclidean action for an instanton solution describing the transition between the two minima of a particle moving in the upside-down potential $V_E = -V(x)$. The factor in front of the exponential can be computed evaluating the fluctuation of the action up to quadratic order around the instanton action. We refer to Appendix 1A for details.
















































\newpage

% Por José Ili\'c García
% Este documento está actualizado a las guías del sistema de bibliotecas de la Universidad de Concepción, solo debes cambiarles los datos correspondientes a tu situación.

\documentclass[spanish, a4paper, 12pt, twoside, openany]{book} 

% a4papper: Tamaño de papel
% 12pt:     Tamaño de fuente 
% twoside:  Escribir capítulos por ambas caras (Solo sirve si se usa Book, de lo contrario no importa)
% openany:  Inicia capítulos en números pares e impares. 
% book:     Formato libro, está dividido por capítulos. Otros formatos: report, letter, article, etc.



% -------------- Setup, no cambiar ---------------
\usepackage{textcomp}                           % Agrega muchos símbolos extra para utilizar en el texto
\usepackage[T1]{fontenc, url}                   % Para los enlaces url y las fuentes
\usepackage[utf8]{inputenc}                     % Introduce la codificación utf8 que permite el uso de tildes y otros.
\usepackage{titlesec}                           % Suplemento a fancyhdr
%\setcounter{secnumdepth}{4}                     % El contador de figuras, secciones, tablas, etc. llega hasta profundidad de 4. i.e. 1.1.1.1
\usepackage{multirow}                           % Extensión para tablas
%\usepackage{minted}                             % Permite darle formato al código que se muestre.
\usepackage{adjustbox}                          % Permite mas versatilidad para cosas como includegraphics.
\usepackage{graphicx}                           % Para poner imágenes
\usepackage{amsmath,amssymb,amsfonts,dsfont,mathrsfs,amsthm,mathtools}          % paquetes matemáticos
%\usepackage{parskip}                            % Elimina la sangría
\urlstyle{sf}                                   % Tipo de url
\usepackage{color}                              % Agrega colores.
\usepackage{cancel}
\usepackage{subcaption}                         % Permite subtextos en las figuras. Es incompatible con el paquete subfig y subfigure
\usepackage[toc,page]{appendix}                 % Permite el uso de apéndices
\usepackage{chngcntr}                           % Necesario para enumerar las tablas



\counterwithin{table}{section}                  % Numeración de tablas 
\counterwithin{figure}{section}                 % Numeración de figuras
\numberwithin{equation}{section}                % Numeración de ecuaciones
\hyphenpenalty=100000                           % No se dividen las palabras al terminar una línea
\sloppy                                         % Ayuda a que las cosas no se salgan de los márgenes.
\raggedbottom                                   % Hace que el tamaño de las páginas llegue hasta donde llegue el texto, no agrega espacio vertical
\usepackage{xparse,nameref}                     % Ayuda en la creación de nuevos comandos
\usepackage[bottom,hang,flushmargin]{footmisc}  % notas al pie de página están fijas al fondo y sin sangría
\interfootnotelinepenalty=10000                 % Previene que las notas al pie de página se pasen a la siguiente página cuando son muy largas
\usepackage{lipsum}                             % Para crear texto de relleno.

% --------- Editar de aquí en adelante --------

% ----- Apariencia e idioma ----- 
\usepackage[spanish]{babel}                                         % Idioma
\graphicspath{{Images/}{../Images/}}                                % Dónde estarán las imágenes
\usepackage[left=2.5cm,top=4cm,bottom=2.5cm,right=2.5cm]{geometry}    % Márgenes del documento
\usepackage{setspace}                                               % Permite elegir el interlineado
\linespread{1.3}                                                    % Interlineado de uno y medio. 1.6 es interlineado doble.
\usepackage{microtype}                                              % Permite la modificación de los caracteres.
\usepackage{ragged2e}

% ----- Secciones ----- % ESTA PARTE SE UTILIZA EN CASO DE USAR LA CLASE ARTICLE

% \titleformat*{\section}{\LARGE\bfseries}                  % Forma del título de \section 
% \titleformat*{\subsection}{\Large\bfseries}               % Forma del  título de \subsection
% \titleformat*{\subsubsection}{\large\bfseries}            % Forma del  título de \subsubsection 

% Las siguientes tres líneas crean el comando \paragraph con la forma del título correcta.

% \titleformat{\paragraph} 
% {\normalfont\normalsize\bfseries}{\theparagraph}{1em}{}
% \titlespacing*{\paragraph}
% {0pt}{3.25ex plus 1ex minus .2ex}{1.5ex plus .2ex}
%-----------------------------------------------

% ----- Figuras y tablas ----- 
\usepackage{fancyhdr}                           % Permite formatear las cabeceras, pies, enumeración, etc.
\usepackage{subfiles}                           % Para agregar los capítulos que se escriben aparte.
\usepackage{array}                              % Para ordenar texto y ecuaciones.
\usepackage[rightcaption]{sidecap}              % Permite agregar texto lateral
\usepackage{wrapfig}                            % Permite poner figuras con texto al rededor.
\usepackage{float}                              % Permite poner figuras en cualquier lugar.
\usepackage[labelfont=bf]{caption}              % Texto en negrita para descripciones (\caption)
\usepackage[para]{threeparttable}               % Tablas vistosas, mirar antes de utilizar.
\usepackage{url}                                % Permite el uso de enlaces URL.
\usepackage[table,xcdraw,dvipsnames]{xcolor}    % Agranda la cantidad de colores.
\usepackage{makecell}                           % Ayuda en la creación de tablas
\usepackage{hhline}                             % Agranda las opciones de las líneas
\usepackage{textcomp}                           % Símbolo de derechos de autor
\usepackage{graphicx}

% ----- Referencias -----
%\usepackage{natbib}
\usepackage[noadjust]{cite}
% Ambiente de referencias utilizado.
\bibliographystyle{unsrt}                                                 % Estilo de referencias APA.
\def\biblio{\clearpage\bibliographystyle{apalike}\bibliography{References}} % Define el comando \biblio para referencias en subarchivos- NO CAMBIAR


% ----- Cabecera y pies -----
\pagestyle{fancy}                           % Se define el estilo fancy
\fancyhead[RO,LE]{\thepage}                 % Número de página en la izquierda para par y derecha para impar
\fancyhead[RE,LO]{\nouppercase{\rightmark}} % Nombre del capítulo en la derecha para par y la izquierda para impar en la cabecera
%\renewcommand{\headrulewidth}{0pt}         % Cambiar para línea más gruesa
\fancyfoot{}                                % Saca el número de la página abajo.

\fancypagestyle{plain}{                     % Se redefine el estilo automático (plain) para que calce con el resto. En particular la 1ra página de cada capítulo
\fancyhf{}                                  % Elimina la cabecera y los pies
\fancyhead[RO,LE]{\thepage}                 % Número de página en la izquierda para par y derecha para impar
\fancyhead[RE,LO]{\nouppercase{\leftmark}}  % Nombre del capítulo en la derecha para par y la izquierda para impar en la cabecera
%\renewcommand{\headrulewidth}{0pt}         % Cambiar para línea más gruesa
\fancyfoot{}                                % Elimina el número de la página abajo
}

%------------------- Cabecera del Resumen y Agradecimientos--------------

\fancypagestyle{resumen}{                   % Se redefine el estilo resumen para que calce con el resto. 
\fancyhf{}                                  % Elimina la cabecera y los pies
\fancyhead[RO,LE]{\thepage}                 % Número de página en la izquierda para par y derecha para impar
\fancyhead[RE,LO]{\nouppercase{Resumen}}    % Nombre del capítulo en la derecha para par y la izquierda para impar en la cabecera
%\renewcommand{\headrulewidth}{0pt}         % Cambiar para línea más gruesa
\fancyfoot{}                                % Elimina el número de la página abajo
}

\fancypagestyle{abstract}{                  % Se redefine el estilo resumen para que calce con el resto. 
\fancyhf{}                                  % Elimina la cabecera y los pies
\fancyhead[RO,LE]{\thepage}                 % Número de página en la izquierda para par y derecha para impar
\fancyhead[RE,LO]{\nouppercase{Abstract}}   % Nombre del capítulo en la derecha para par y la izquierda para impar en la cabecera
%\renewcommand{\headrulewidth}{0pt}         % Cambiar para línea más gruesa
\fancyfoot{}                                % Elimina el número de la página abajo
}

\fancypagestyle{agradecimientos}{                   % Se redefine el estilo resumen para que calce con el resto. 
\fancyhf{}                                          % Elimina la cabecera y los pies
\fancyhead[RO,LE]{\thepage}                         % Número de página en la izquierda para par y derecha para impar
\fancyhead[RE,LO]{\nouppercase{Agradecimientos}}    % Nombre del capítulo en la derecha para par y la izquierda para impar en la cabecera
%\renewcommand{\headrulewidth}{0pt}                 % Cambiar para línea más gruesa
\fancyfoot{}                                        % Elimina el número de la página abajo
}

% ----- Cabecera de la portada ----- 
\fancypagestyle{frontpage}{             % Se define el estilo frontpage.
	\fancyhf{}                          % Elimina la cabecera y los pies
	\renewcommand{\headrulewidth}{0pt}  % Elimina líneas en cabecera
	\renewcommand{\footrulewidth}{0pt}  % Elimina líneas en pies
	\vspace*{1\baselineskip}
	
	\   % Línea 2 de la cabecera derecha, debe ir el nombre de la facultad en mayúsculas.
	\fancyhead{\centering \includegraphics[width=1.8in]{Images/logo-unap2.jpg}                       % Línea 1 de la cabecera derecha, debe ir el nombre de la UdeC en mayúsculas.
	\linebreak\linebreak       Instituto de Ciencias Exactas y Naturales \\ Facultad de Ciencias} % LOGO que irá en la cabecera izquierda.
}

% ----- Enlaces clickeables --------
\usepackage{hyperref}   % Permite que todo el documento sea clickeable.
\newcommand\myshade{85} % Permite la redefinición de colores a gusto del usuario

% Para elegir colores propios mirar los nombres relacionados con dvipsnames, aquí un url con los nombres de dvipsnames: https://www.overleaf.com/learn/latex/Using_colours_in_LaTeX

\colorlet{mylinkcolor}{DarkOrchid}   % Define Color de los enlaces clickeables internos (índice, referencias cruzadas, etc.). En este caso DarkOrchid
\colorlet{mycitecolor}{YellowOrange} % Define Color de las citas ubicadas en las referencias. En este caso YellowOrange
\colorlet{myurlcolor}{Aquamarine}    % Define Color de los enlaces url. En este caso Aquamarine

% Para dejar el documento sin texto en colores cambiar las tres líneas anteriores a Black.

\hypersetup{  %Define la forma en que se verán las cosas clickeables.
  	linkcolor  = mylinkcolor!\myshade!black,    % Aplica el color definido arriba. En este caso DarkOrchid
  	citecolor  = mycitecolor!\myshade!black,    % Aplica el color definido arriba. En este caso YellowOrange
  	urlcolor   = myurlcolor!\myshade!black,     % Aplica el color definido arriba. En este caso Aquamarine
  	colorlinks = true,                          % Elimina las cajas al rededor de lo clickeable y lo reemplaza por palabras a color.
}


%--------------------------------------------------------------------------------------------------------------------------
% ------------------------------------------ Aquí empieza el documento ----------------------------------------------------
%--------------------------------------------------------------------------------------------------------------------------

\newcommand*{\diag}{\operatorname{diag}}
\newcommand*{\Tr}{\operatorname{Tr}}
\newcommand*{\tr}{\operatorname{tr}}
\newcommand{\diff}[1]{\text{d}#1}
\newcommand{\Diff}[1]{\text{D}#1}
\newcommand{\Lie}{\mathcal{L}}
\newcommand{\Lag}{\mathscr{L}}

\usepackage{hyperref}

\begin{document}
\def\biblio{}   % Resetea el comando biblio, de lo contrario una lista de referencias será producida después de cada capítulo
                % resets the biblio command, if not here a new reference list will be produced after every chapter


\begin{titlepage}
	
	\newgeometry{top=1 in, bottom=1 in, left=1 in, right= 1 in} 
	
	\thispagestyle{frontpage}
	
	\begin{center}
		
		\vspace*{8\baselineskip}
	
		
		{\Huge \textbf{Renormalización conforme de teorías escalar-tensor de la gravedad\\}}%No abreviar, no subrayar, no usar comillas. Se escribe completamente en mayúsculas
		        \vspace*{1.5\baselineskip}

		% \large{\textit{subtítulo}}\\ %No abreviar, no subrayar, no usar comillas. Solo la priemra letra usa mayúsculas
		
        \vspace*{1\baselineskip}

		\large{\textbf{Por: Mairym Eliana Busnego Barrientos}}\\ %Nombre como aparece en registro académico
		
		\vspace{1,5\baselineskip}
		
		\large{Tesis presentada a la Facultad de Ciencias de la Universidad Arturo Prat y al Instituto de Ciencias Exactas y Naturales y  para optar al grado académico de Magíster en Ciencias con Mención en Física Teórica.} %Cambiar a grado académico correspondiente. 
		
		\vspace{1,5\baselineskip}
		\begin{center}
            Marzo 2023 \\
            Iquique, Chile
            \end{center}
\vspace{1.5\baselineskip}

		\large{\textbf{Directores de tesis: Dr. Cristóbal Corral y  Dr. Nelson Merino}}\\ %Nombre del profesor que dirigió el trabajo, de la comisión informante y otros asesores precedidos por los títulos correctos de la unidad académica. Si hay página de calificaciones, solo se menciona al profesor guía en la portada
            \large{\textbf{Comisión: }}\\
	\end{center}
	
	\vspace*{4\baselineskip}
	
	% La cabecera de esta página se cambia en el documento principal, en "cabecera portada"
	
\end{titlepage}   % Genera la portada desde el archivo Portada.tex . Para editar este archivo ir a Otros -> Portada.tex
\vfill

%----------------Página de derechos de autor: elegir entre a) o b) y borrar/comentar la opción NO utilizada-----------------
\thispagestyle{empty}
\mbox{}                         % Ayuda a bajar el texto
\vfill                          % Deja el texto al fondo
\textcopyright\ 2024, Borja Diez Bravo \\ % Derechos de autor
%a)
Ninguna parte de esta tesis puede reproducirse o transmitirse bajo ninguna forma o por ningún medio o procedimiento, sin permiso por escrito del autor.\\\\
%b)
Se autoriza la reproducción total o parcial, con fines académicos, por cualquier medio o procedimiento, incluyendo la cita bibliográfica del documento
\vspace{1cm}    % lo separa del fondo
\restoregeometry % Devuelve los márgenes después de la portada


%----------------Página de calificaciones (opcional), descomentar para generar-----------------

% Editar en Otros -> Calificaciones.tex

%%Las calificaciones entregan la misma información de la portada y se le agrega la nota y la firma. Esta página es opcional.
\begin{titlepage}
	
	\newgeometry{top=1 in, bottom=1 in, left=1 in, right= 1 in} 
	
	\thispagestyle{frontpage}
	
	\begin{center}
		
		\vspace*{4\baselineskip}
	
		
		{\Huge \textbf{TÍTULO PRINCIPAL\\}}%No abreviar, no subrayar, no usar comillas. Se escribe completamente en mayúsculas
		        \vspace*{1.5\baselineskip}

		\large{\textit{subtítulo}}\\ %No abreviar, no subrayar, no usar comillas. Solo la priemra letra usa mayúsculas
		
        \vspace*{1,5\baselineskip}

		\large{\textbf{Por: Autor}}\\ %Nombre como aparece en registro académico
		
		\vspace{1,5\baselineskip}
		
		\large{Tesis presentada a la Facultad de Ciencias Físicas y Matemáticas de la Univerisdad de Concepción para optar al grado académico de Magister en Ciencias con Mención en Física.} %Cambiar a grado académico correspondiente. 
		
		\vspace{1,5\baselineskip}
		Diciembre 2019\\ %Mes y año de la tesis, solo primera letra del mes en mayúsculas
		Concepción, Chile %Ciudad y país de publicación.
\vspace{1,5\baselineskip}

		\large{\textbf{Profesor Guía: Nombre}}\\ %Nombre del profesor que dirigió el trabajo, de la comisión informante y otros asesores precedidos por los títulos correctos de la unidad académica.  
		%Al margen derecho del profesor guía se estampa la firma y/o calificación.
		

	\end{center}
	

	
\end{titlepage}         % Genera la pagina de calificaciones del archivo calificaciones.tex
%\restoregeometry                           % Devuelve los márgenes después de la página

%\pagenumbering{gobble}         % Suprime la numeración de páginas
%\thispagestyle{plain}          % suprime el encabezado
%\clearpage\mbox{}\clearpage    % Agrega página en blanco

%----------------Página de dedicatoria (opcional), descomentar para generar ---------------------------------

% Editar en Otros -> Dedicatoria.tex

%La dedicatoria debe ser ubicada en el extremo inferior derecho de una página vacía, está página es opcional
\thispagestyle{empty}
\begin{flushright}
\vfill
\textit{Now the world can be an unfair place at times\\
But your lows will have their compliment of highs\\
And if anyone should cheat you\\
Take advantage of or beat you raise your head\\
And wear your wounds with pride.}
\vfill
\end{flushright}%Escribir aquí la dedicatoria
   % Genera la pagina de calificaciones del archivo calificaciones.tex
\restoregeometry        % Devuelve los márgenes después de la página


%-----------------Página de agradecimientos (opcional), se incluye normalmente-------------------

% Editar en Otros -> Agradecimientos.

\pagenumbering{roman}                            % Empieza la enumeración romana en minúsculas, para mayúsculas usar Roman.
\newpage
\addcontentsline{toc}{chapter}{AGRADECIMIENTOS}  % Agrega esta sección al índice
\section*{AGRADECIMIENTOS}                       % Debe ir en mayúsculas por reglamento de la UdeC, tiene asterisco para no ser numerada.
    \subfile{Otros/Agradecimientos}              % Llama al archivo Agradecimientos en la carpeta Otros.


%-----------------Página de resumen (abstract)-------------

% Si la unidad académica lo requiere, se edita en  Otros -> Resumen.tex . El mismo resumen puede ser incluido en inglés (abstract) en la página siguiente, para agregarlo hay un espacio destinado en el mismo archivo antes mencionado.

\newpage
\addcontentsline{toc}{chapter}{Resumen} % Agrega esta sección al índice
\section*{Resumen}                      % Con asterisco para que no sea numerada.
    \subfile{Otros/Resumen}             % Llama al archivo Resumen en la carpeta Otros.

%--------------Página de índice.  

%\nocite{*}     % Des-comentar si se desea que TODAS las referencias sean impresas en la lista de referencias, incluyendo las que no fueron finalmente citadas en el texto.

\newpage
{\setstretch{1.0}   % Interlineado de la lista.
\tableofcontents
}

% \newpage
% {\setstretch{1.0} 
% \listoftables}

% \newpage
% {\setstretch{1.0} 
% \listoffigures}


\newpage
\addtocontents{toc}{\protect\setcounter{tocdepth}{4}}   % La profundidad del índice queda en 4, 1.1.1.1
\pagenumbering{arabic}                                  % Comienza la numeración arábiga (números normales)
\setcounter{page}{1}                                    % Comienza el contador de páginas en 1

% A continuación se dejan nombres de diversos capítulos o secciones, para cambiar el nombre del archivo tan solo se debe hacer en la carpeta "capitulos" y luego llamarlos de la forma correcta en "\subfile{Capitulos/nuevonombre}".
% Los nombres de los archivos no pueden llevar tíldes ni espacios para el correcto funcionamiento del compilador, esto no tiene nada que ver con que tengan o no tilde en el documento final.


blablabla


\newpage
\renewcommand\refname{Referencias}          % Nombre para la lista de referencias, también se utiliza "Bibliografía"
{\setstretch{1.0}                           % Interlineado de las referencias 
\addcontentsline{toc}{chapter}{Referencias} % Cambia el nombre de la lista de referencias en el índice 
\bibliography{Referencias.bib}              % Agrega las referencias al documento, estas se ubican en el archivo Referencias.bib
}

\newpage
\renewcommand{\appendixpagename}{Apéndices}     % Nombre al inicio.
\addcontentsline{toc}{chapter}{Apéndices}       % Agrega "Apéndices" al índice

% \appendix   % Empieza el ambiente de apéndices, desde ahora en adelante los capítulos, secciones, tablas, figuras, etc. vuelven a empezar su numeración

% \chapter{Test}                      % Nombre del apéndice, el capítulo será "Apéndice A" en vez de "Capítulo 1"
%     \subfile{Capitulos/Apendice}    % Llama al archivo apéndice
% \clearpage

% Este segundo apéndice tiene información general sobre LaTeX. Está comentado para que no aparezca en tu documento, pero si lo deseas puedes descomentarlo para ver su contenido.

%\chapter{Algunos consejos sobre \LaTeX{}}
%    \subfile{Otros/Consejos}
%\clearpage

% ÉXITO EN TU TESIS

\end{document} 




% Bibliography

%% [A] Recommended: using JHEP.bst file
%% \bibliographystyle{JHEP}
%% \bibliography{biblio.bib}

%% or
%% [B] Manual formatting (see below)
%% (i) We suggest to always provide author, title and journal data or doi:
%% in short all the informations that clearly identify a document.
%% (ii) please avoid comments such as "For a review'', "For some examples",
%% "and references therein" or move them in the text. In general, please leave only references in the bibliography and move all
%% accessory text in footnotes.
%% (iii) Also, please have only one work for each \bibitem.



\newpage
\bibliographystyle{JHEP}
\bibliography{biblio.bib}
\end{document}
